% Gerado por scripts/gerar_tabelas.py --- não editar à mão.
\begin{table}[htbp]
    \centering
    \caption{Decomposição de Oaxaca-Blinder do hiato salarial}
    \label{tab:hiato}
    \footnotesize
    \renewcommand{\arraystretch}{1.25}
    \begin{tabular}{llrrr}
        \toprule
        \rowcolor{black!8}
        \textbf{Comparação} & \textbf{Ano} & \textbf{Hiato bruto} & \textbf{Parcela explicada (log)} & \textbf{Parcela não explicada (log)} \\
        \midrule
        Branca vs Parda & 2007 & 64,0\% & 0,388 & 0,107 \\
         & 2015 & 42,0\% & 0,274 & 0,077 \\
         & 2025 & 38,1\% & 0,237 & 0,086 \\
        \cmidrule{1-5}
        Branca vs Preta & 2007 & 51,1\% & 0,262 & 0,151 \\
         & 2015 & 41,1\% & 0,194 & 0,151 \\
         & 2025 & 30,9\% & 0,145 & 0,124 \\
        \cmidrule{1-5}
        Homem vs Mulher & 2007 & $-$0,9\% & $-$0,150 & 0,141 \\
         & 2015 & $-$0,7\% & $-$0,132 & 0,125 \\
         & 2025 & 7,8\% & $-$0,058 & 0,134 \\
        \bottomrule
    \end{tabular}
    \\[4pt]
    Parcelas em logaritmo da remuneração; referência de Neumark. Hiato bruto $= e^{E+N} - 1$, com $E$ e $N$ as parcelas explicada e não explicada. Fonte: Elaborado pelo autor com dados da RAIS \cite{RAIS2024}.
\end{table}
